% !TEX root = begin.tex
% !TeX spellcheck = nl_NL
% \titelpagina definieert de titlepage
% 1 = Naam van het plaatje in directory ./fig. 
%     Neem een hoog resolutie plaatje.
%     De hoogte moet ongeveer 0.5 x de breedte zijn.
%>titelpagina
%\titelpagina{logo-hr.png}
%<titelpagina
% commando voor versiehistory tabel
% #1 = tabel inhoud
%      tabel heeft 4 kolommen:
%      Datum & Versie & Omschrijving & Auteur
%      Gebruik \\ \line om de rijen te scheiden.
%>versiehistorie
\versiehistorie{
	7-4-2026 & \huidigeversie &
	Biblografie toegvoegt en verweizingen naar bronnen in de tekst & JHT
	\\ \hline
	5-03-2026 & 1.1 &
	tabellen toegevoegt en code opmaak verbeterd & JHT
	\\ \hline
	24-03-2026 & 1.0 &
	Driefasen toegevoegt & JHT
	\\ \hline
	24-03-2026 & 0.2 &
	Hoofdstukken 1, 2 en 3 toegevoegt
	& JHT
	\\ \hline
    24-03-2026 & 0.1 & 
	Instellingen, versiehistory en titelpagina toegevoegt 
	& JHT 
}

%<versiehistorie
%voorbeeld gebruik footnote
% \footnotetext{Toelichting versiecodering $A$.$Bc$: $A$ = grote aanpassing, $B$ = kleine aanpassing, $c$ = taal- of wiskundige
% correcties.}

%>licentie    
\vspace*{\stretch{100}}
\par